\chapter{Tubular Agents}
\label{ch:tub}
% ===================================================================

This part built learners at four sizes and then ran the architecture downward to
a synapse. It has one thing left to do, and it is a thing that can only be done
at the end, because it requires all of the machinery at once.

Return to \chSci and to the distinction it drew in \S\ref{sci:sec:trophic}. A
\emph{source} is a persistent emitter: rate-limited, non-exclusive,
non-adaptive. \emph{Prey} is a linear send: exhaustible, exclusive, adaptive.
Same sort, same rewrite rule, different multiplicity. That distinction was put
to work immediately --- it is what Corollary~\ref{sci:cor:intelligence} rests on,
and it is why plants do not need brains --- and then it was set aside while the
unit of learning moved four times.

We now have what \chSci did not: an account of what an individual is
(Definition~\ref{eng:def:indiv}), of which of its channels face outward
(Definition~\ref{eng:def:three}), of what ownership of a pathway buys
(Corollary~\ref{eng:cor:owndepth}), and of how one level's outside becomes the
next level's inside (Proposition~\ref{eng:prop:tower}). With those in hand the
trophic distinction says more than it did. It says something about shape.

The claim of this chapter is that an animal which must open its food is,
topologically, a tube; that this is a consequence of an access profile rather
than an observation imported from anatomy; and that the condition under which it
holds can be stated as an inequality between two quantities the framework
already meters. That last part is what makes the chapter worth ending a part
with rather than a remark worth making in passing. A criterion that says when a
tube is forced also says when one is not, and a claim that predicts bags as well
as tubes has a test set rather than an illustration.

Much of what follows is conjecture and is labeled as such. The chapter is the
part's closing gesture rather than its last proof. But the gesture points at
something checkable, and the checking has begun.

\section{The observation}
\label{tub:sec:obs}

The reason for eating is to gain access to the token supply held inside the
computation being consumed.

There is an advantage to minimizing the time spent \emph{exclusively} on
extraction. If some part of the extraction can be carried on internally --- so
that the consumer can meanwhile forage, mate, form hypotheses, and run assays
--- then extraction proceeds concurrently with everything else the learner does,
and by the accounting of \S\ref{sci:sec:game} that concurrency is worth paying
for.

Eating therefore has two phases. There is a gross phase, which gets the relevant
material across the boundary; and an internal phase, in which the remaining
separation runs while the learner is engaged elsewhere.

The internal phase leaves a residue: the part that was not token. It must leave.
Material in, refuse out.

\begin{observation}[Refuse is not an extra postulate]
\label{tub:obs:refuse}
The residue exists because prey are computations and sources are not. A source
emits tokens and nothing else; breaking prey open yields token entangled with
code. The non-token part of a structured prey term is the refuse, and nothing
beyond Table~\ref{sci:tab:trophic} of \chSci is required to produce it.
\end{observation}

\section{The gut is an internalized enzyme}
\label{tub:sec:ident}

\chEng put a converting medium between two learners and ranked media by behavior
in Table~\ref{eng:tab:ladder}: relay, lossy, corrupting, reading, converting,
with the enzyme of Definition~\ref{eng:def:engine} at the top rung. The
mycorrhizal network was the external instance of that profile.

\begin{observation}[The gut is the enzyme profile drawn inside]
\label{tub:obs:gut}
The digestive apparatus is a converting medium held within the individual's
boundary namespace: a chain of interposed channels running from an intake
channel to an expulsion channel, with the whole chain owned.
\end{observation}

This is not an analogy between two systems. It is one access profile appearing
at two positions relative to a boundary, which is exactly what
Proposition~\ref{eng:prop:tower} licenses. The observation that makes the
topology non-metaphorical is the next one, and it is a definitional consequence
rather than a new claim.

\begin{observation}[The lumen is outside]
\label{tub:obs:lumen}
The interior of the tube is not in $\Chn_{\mathrm{in}}(B)$. It is
$\Chn_{\partial}(B)$: owned, but open.
\end{observation}

A tube is therefore the construction that surrounds a region of the outside with
the inside. It is the tower's generating step applied for a specific purpose ---
to acquire boundary surface without acquiring exposure, in the sense of the
exposure lemma of \S\ref{sci:sec:exposure}. The biological commonplace that the
gut lumen is topologically external \cite{stevenshume1995} is, on this reading,
the same fact as the framework's fractality, and not a coincidence to be noted
alongside it.

\section{Organs are communities}
\label{tub:sec:organs}

Here is where the part's four moves pay off, and where the picture is not the
one the trophic section of \chSci could have drawn.

\chComp made a composite out of two learners sharing a namespace, and \chEng
gave a criterion for where one individual ends and the next begins. Put those
together with Observation~\ref{tub:obs:lumen} and something follows that is
worth stating slowly. The boundary namespace of an individual is owned but open.
Owned means it can be let. Open means whatever occupies it is not inside.

\begin{observation}[The boundary namespace can be let]
\label{tub:obs:let}
A population of learners occupying $\Chn_{\partial}(B)$ is neither part of $B$'s
internal computation nor foreign to it. It performs conversions on $B$'s intake
under $B$'s ownership, and by Observation~\ref{tub:obs:lumen} it does so
outside.
\end{observation}

A digestive organ, on this reading, is not a specialized tissue that happens to
convert. It is a community of learners tenanted in a boundary namespace, and the
gradient of conversions along a gut --- one flavor handed to the next --- is a
succession in the sense of \S\ref{sci:sec:trophic}, running in space rather than
in time.

Two consequences make this more than a re-description.

The first is that it removes an objection. \chEng noted that a market inside an
individual is Coase's theory of the firm \cite{coase} arriving unbidden, and that
the arrival is awkward: if conversion is cheaper inside than across a boundary,
why is there a boundary? Under Observation~\ref{tub:obs:let} the awkwardness
dissolves, because the market is not inside. The rented enzymes sit in
$\Chn_{\partial}$, where pricing is still across a cut. The firm has an interior
in which no market operates and a lumen in which one does, and the two are
different namespaces.

The second is that it makes the herbivore's symbionts the general case rather
than a curiosity. By Theorem~\ref{eng:thm:epistemic}, tokens held in a foreign
flavor cannot pass without an enzyme linking the flavors. A learner that eats far
from itself must therefore either manufacture the conversions or rent them, and
renting is what a boundary namespace is for. The rumen, the hindgut, and the
microbiota that occupy them are not an accessory to digestion
\cite{karasov2013}; they are digestion, performed by a population that the host
does not own but whose namespace it does.

Whether the letting can be done \emph{without} the enclosure is a question this
section cannot yet ask. It is the question \S\ref{tub:sec:testset} turns out to
need, and Observation~\ref{tub:obs:evert} is the answer.

\section{What follows, and what would have to be shown}
\label{tub:sec:props}

Two of the following are argued here. The rest are conjectures, and are labeled
as such because the part's standing commitment is that a claim be stated
precisely enough to be wrong.

\begin{proposition}[Two ends, not one]
\label{tub:prop:twoends}
Intake and expulsion cannot share a name without a mutual exclusion. On a single
shared channel an offer of refuse and a receipt for intake form a redex: the
organism re-ingests its residue, or the residue blocks the intake. A
single-opening organism is therefore not impossible, but its intake and
expulsion must alternate on the shared name, and its throughput is bounded by
that alternation.
\end{proposition}

This is the argument of \chGame one level up. There, a line signaling a threat on
a square's move channel would have rendezvoused with the square's receipt and
played the move; perception must not be able to act, and the separation was
enforced by the choice of namespace. Here, egestion must not be able to feed, and
the separation is enforced the same way.

What Proposition~\ref{tub:prop:twoends} does not say is when the alternation is
cheap enough to live with. That is the business of \S\ref{tub:sec:rate}, and it
is where the chapter acquires its test set.

\begin{proposition}[Non-phagotrophs are not tubes]
\label{tub:prop:notube}
Where an agent admits no bulk material with a non-absorbable fraction, both of
the tube's ends lose their reason: there is nothing whose acquisition must be
exclusive, and by Observation~\ref{tub:obs:refuse} there is no residue to expel.
Source access is one way to be in that position. It is not the only way.
\end{proposition}

An earlier statement of this argument said \emph{autotrophs} where
Proposition~\ref{tub:prop:notube} says non-phagotrophs, and the correction is not
cosmetic. \S\ref{tub:sec:testset} is where it is forced, and it is the most
useful thing the test set did.

In the source case, Proposition~\ref{tub:prop:notube} and
Corollary~\ref{sci:cor:intelligence} are the same sentence about the same access
profile, differing only in what is being purchased. One says that a
source-feeder need not buy continuing inquiry. The other says it need not buy a
topology. Plants do not need brains and plants are not tubes for one reason, not
two. And the prediction is not merely that a plant needs no tube but that it
should buy boundary area directly instead: a leaf and a root system are that
purchase.

\begin{conjecture}[Separation is metered but token-neutral]
\label{tub:conj:sep}
In the harvest of a structured prey term, the rewrite steps separating token
from residue cost and do not pay.
\end{conjecture}

\begin{conjecture}[The serialization penalty]
\label{tub:conj:serial}
If separation is sequential with the learner's control loop, the learner pays
its burn rate across the whole separation with no inquiry, no foraging and no
mating during it. Writing Proposition~\ref{sci:prop:foraging} with a sequential
fraction $\sigma$ should yield an Amdahl-shaped bound \cite{amdahl1967} on
sustainable harvest rate, so that the advantage of concurrency emerges as a
strict inequality with a threshold rather than as a preference.
\end{conjecture}

\begin{conjecture}[Concurrency forces internalization]
\label{tub:conj:intern}
Concurrency as such is free in \rhoc{}: it is parallel composition. The content
is therefore not that extraction \emph{can} be concurrent but \emph{where the
extraction channels live}. If separation runs on unowned boundary channels,
other learners may rendezvous with partially processed material, and exclusivity
of a partly extracted harvest is lost. Exclusivity requires the separation chain
to be owned, which is Corollary~\ref{eng:cor:owndepth} specialized to digestion.
\end{conjecture}

Conjecture~\ref{tub:conj:intern} is the step that converts a remark about
scheduling into a claim about topology, and it is the one we would most like to
see settled.

\begin{conjecture}[Transient versus persistent tubes]
\label{tub:conj:transient}
A food vacuole is a tube built and dissolved per meal; a gut is one maintained.
Persistence carries a maintenance cost against the metabolic stack, so there
should be a threshold harvest rate above which the permanent tube beats the
transient one --- amoeba below it, annelid above.
\end{conjecture}

Proposition~\ref{sci:prop:composite} already says the overlap region between
strategies is populated by composites rather than atoms, so a graded answer here
is expected rather than disappointing.

\begin{conjecture}[Gut depth tracks incommensurability]
\label{tub:conj:depth}
The number of interposed converting stages should scale with the flavor distance
between prey namespace and consumer namespace, in the sense of
\S\ref{eng:sec:incomm}.
\end{conjecture}

The comparative anatomy is corroborating and instructively imperfect. Carnivores
eat close relatives and have short, simple guts; herbivores eat far ones and
have long ones with large fermentation chambers. But the association resisted
statistical demonstration for a long time, and when it was finally shown across
$519$ mammal species the effect turned out to sit in the large intestine rather
than the small, to be inseparable from phylogeny, and to admit enough scatter
that the authors decline to call it a fixed law \cite{duquecorrea2021}. A
framework predicting a monotone relation between flavor distance and stage count
must account for that scatter, and the honest reading is that different
morphological solutions exist for the same conversion problem --- which is what
Proposition~\ref{sci:prop:composite} would lead one to expect.

\begin{conjecture}[The second heterotroph's expense]
\label{tub:conj:second}
Corollary~\ref{sci:cor:intelligence} derives ``intelligence is a heterotroph's
expense'' from Proposition~\ref{sci:prop:limit} together with
Remark~\ref{sci:rem:structural}. The tube should be a second corollary of the
same proposition: exclusivity together with structure forces both purchases.
\end{conjecture}

If Conjecture~\ref{tub:conj:second} goes through it is the strongest result
available here, because it makes a gross anatomical fact a consequence of an
access profile. Intelligence and the tube become siblings rather than neighbors.

\section{Rate of exchange: why breathing is not eating}
\label{tub:sec:rate}

Respiration is also an exchange across a boundary. Something is admitted, a
fraction of it is extracted, and the remains are expelled. Yet the mammalian
lung is a bag and not a tube.

The obvious thought is that the discriminating variable is \emph{rate}:
digestion runs orders of magnitude slower than the respiratory cycle, and it is
the slowness that forces the second opening. That thought is pointing at
something real. Rate as such is nevertheless the wrong variable, and saying why
sharpens the chapter rather than weakening it.

\subsection{Respiration goes tube where one would not expect it}

Ventilation is unidirectional in fish --- water in at the mouth, out at the
opercular openings --- and unidirectional in birds, where air moves the same way
through the parabronchi during both phases of the cycle. Unidirectional
pulmonary flow is not confined to birds. It has been demonstrated in alligators
\cite{farmersanders2010} and in the savannah monitor \cite{schachner2014}, and
Farmer's review argues from that distribution that the trait is ancestral for
diapsids and therefore \emph{not} an adaptation for high rates of gas exchange,
since the animals that have it include ectotherms that do not fly; the candidate
functions offered instead are reducing the work of breathing, evaporative water
loss, and heat loss \cite{farmer2015}.

So a respiratory system can be a tube while running fast, and a bag while
running fast. High rate is neither necessary nor sufficient. What varies across
those cases is the ratio of demand to what a tidal exchanger can supply ---
water holds a small fraction of air's oxygen and is far more viscous --- rather
than the rate itself.

\subsection{Two parameters, not one}

\begin{definition}[Occupancy ratio]
\label{tub:def:lambda}
$\Lambda = \tau_{\mathrm{proc}} / \tau_{\mathrm{enc}}$, where
$\tau_{\mathrm{proc}}$ is the residence time of an admitted load inside the
exchanger and $\tau_{\mathrm{enc}}$ is the interval at which a further load
could profitably be admitted.
\end{definition}

\begin{definition}[Separability]
\label{tub:def:mu}
$\mu \in [0,1]$ measures the degree to which residue can leave past incoming
material at the same aperture without bulk transport. Take $\mu = 1$ for a
residue miscible with the intake and free to leave by diffusion, and $\mu = 0$
for a residue that must be conveyed as a bolus.
\end{definition}

$\Lambda$ is what sets the sequential fraction $\sigma$ of
Conjecture~\ref{tub:conj:serial}. If the control loop is blocked for the whole
residence then $\sigma = \Lambda / (1 + \Lambda)$, and the Amdahl-shaped bound
bites only as $\Lambda$ approaches and passes unity. This is the sense in which
the rate intuition is right: slowness matters, but only relative to the
opportunity structure. A process that takes a day is cheap to serialize if the
next load cannot arrive for a week.

\subsection{Breathing fails on separability, not on rate}

The respiratory residue fraction is close to one --- almost all of the mass
admitted is expelled again --- and bag topology survives that because
$\mu \approx 1$. The mammalian lung in fact tolerates gross re-mixing. An
inspired bolus on the order of a quarter of a liter mixes with something like a
liter and a half of residual gas, six parts spent medium to one part fresh,
every cycle \cite{farmer2015}. It works because the residue is the same phase as
the resource and separation happens by diffusion at the membrane rather than by
transport through the aperture.

Digestion has $\mu \approx 0$. A bolus of partly extracted material cannot
diffuse past incoming food, and the cost of re-mixing is not a few percent of
exchange efficiency but contamination of the batch.

\begin{proposition}[The tube criterion]
\label{tub:prop:criterion}
A tube is forced when the agent is \emph{phagotrophic} --- admitting bulk
material with a non-absorbable fraction --- and $\Lambda \gtrsim 1$ and
$\mu \approx 0$. Failure of any one of the three is sufficient for a bag.
\end{proposition}

Proposition~\ref{tub:prop:criterion} subsumes
Proposition~\ref{tub:prop:twoends} rather than replacing it. The mutual
exclusion on a shared name is what the agent pays; $\Lambda$ is how much of the
cycle that exclusion consumes; $\mu$ is whether the exclusion can be evaded by
letting the two flows share the aperture after all.

The criterion is worth having because it has three clauses and each can fail
separately. A theory that only ever predicts tubes explains nothing. This one
predicts a bag whenever the agent is an osmotroph, or whenever loads arrive
rarely relative to how long they take to process, or whenever the residue is the
same stuff as the intake. Those are three different kinds of animal, and the next
section asks whether they are the right three.

\section{The test set}
\label{tub:sec:testset}

M.~Stay proposed a list of animals with bag-like rather than tube-like digestive
systems --- cnidarians, ctenophores, flatworms, xenacoelomorphs and ophiuroids
--- together with two cases from outside the animals: fungi, which are
heterotrophs, and plants such as \emph{Monotropa uniflora}, which parasitize
fungi rather than photosynthesize. It is the right test set, and it separates
into three kinds of case: one that is factually mistaken, four that the
occupancy ratio predicts, and two that force a correction of scope.

\begin{table}[t]
\centering\small
\fitwidth{%
\begin{tabular}{@{}llp{6.2cm}@{}}
\toprule
\textbf{Case} & \textbf{Verdict} & \textbf{Why} \\
\midrule
Ctenophora        & fails     & has a through-gut \\
Cnidaria          & predicted & low $\Lambda$; duty-cycles when $\Lambda$ rises \\
Platyhelminthes   & predicted & buys area; gutless where osmotrophic \\
Xenacoelomorpha   & predicted & no lumen at all; diffusion-scale bodies \\
Ophiuroidea       & predicted & secondary loss, microphagous diets \\
Fungi             & scope     & osmotroph, not phagotroph \\
Mycoheterotrophs  & scope     & osmotroph, not phagotroph \\
\bottomrule
\end{tabular}}
\caption{Seven proposed counterexamples to the tube argument, set against
Proposition~\ref{tub:prop:criterion}.}
\label{tub:tab:cases}
\end{table}

\subsection{Ctenophora: the counterexample that fails}

Comb jellies are not bags. Time-lapse imaging showed that the ctenophore gut is
unidirectional and functionally tripartite, with waste expelled through terminal
anal pores that are specialized to control outflow, resolving a long-standing
misreading of those pores that had supported the blind-gut picture
\cite{presnell2016,giribet2016}.

This is more than the removal of a counterexample. Ctenophores branch very deep
in the metazoan tree, so a through-gut there is evidence that the construction is
not an inheritance of the bilaterian body plan but something arrived at from an
access profile --- which is this chapter's thesis, and the strongest single piece
of support it has.

\subsection{Cnidaria: the predicted duty cycle}

Mean digestion time in \emph{Aurelia aurita} has been measured at about one hour
\cite{ishii2001}. Against the encounter interval of a drifting tentacle feeder on
dilute plankton that is a low $\Lambda$, and
Proposition~\ref{tub:prop:criterion} accordingly predicts a bag.

The prediction with teeth is what should happen when $\Lambda$ rises, and it has
been observed. In \emph{Pelagia noctiluca} ephyrae feeding on a dense and pulsed
prey field, the gut saturates in about fifteen minutes while digestion takes
about eighteen hours, and the authors read the mismatch as implying a diel
feeding periodicity \cite{gordoa2013}. Fast fill, slow clear, forced
alternation: that is Proposition~\ref{tub:prop:twoends} observed in a single
organism, with the ratio measured.

\subsection{Platyhelminthes and Xenacoelomorpha: area instead of length}

Acoels have no gut lumen at all. The mouth opens into a syncytial digestive
parenchyma without epithelial lining, in which particles are handled by
phagocytosis. At body scales where the diffusion distance is a few cell
diameters there is no bulk transport to arrange and $\mu$ has no work to do. The
branched blind gut of a triclad is the same move at larger scale: it buys
boundary area rather than a second opening, which is the purchase
Proposition~\ref{tub:prop:notube} assigns to the non-phagotroph.

Cestodes are the sharper case. Tapeworms have no digestive system whatever and
take up small organic monomers across the tegument. A cestode is an animal, a
bilaterian, a heterotroph, and descended from ancestors with a through-gut ---
and it is not a tube, because something else did the digestion. No formulation in
terms of heterotrophy can accommodate that. Proposition~\ref{tub:prop:criterion}
accommodates it by the phagotrophy clause, and this is the case that forced the
clause.

\subsection{Ophiuroidea: secondary loss under a diet shift}

Brittle stars lack an anus, as do paxillosid asteroids, and the review
literature catalogues repeated independent losses of the anal opening across
Bilateria, of which these are two \cite{hejnol2015,jangoux1982}. The direction of
change matters: this is loss from a through-gut ancestor rather than a lineage
that never built one, and it is concentrated in microphagous habits ---
suspension and deposit feeding on small particles, where the indigestible
fraction per load is small and $\Lambda$ correspondingly low.

That yields a within-clade test, which is the kind worth wanting. Macrophagous
ophiuroids --- scavengers and predators taking large items --- should show
markedly harder duty-cycling than their microphagous congeners, and if they do
not, Proposition~\ref{tub:prop:criterion} is in trouble on ground of its own
choosing.

\subsection{Fungi and mycoheterotrophs: the scope was stated too widely}

These are not counterexamples. They are a correction, and a useful one, because
they identify the clause the argument had left implicit.

Heterotrophy is a claim about carbon source. The tube argument is a claim about
\emph{phagotrophy}: about admitting bulk material that contains something which
will not be absorbed. A fungus secretes its enzymes outward and takes up
monomers; the residue never crosses the boundary at all. Mycoheterotrophic
plants do the same at one further remove, drawing fixed carbon through a
mycorrhizal interface \cite{leake1994,bidartondo2005}, and \emph{Monotropa
uniflora} is the resolutive case, taking everything it has through associations
with Russulaceae.

\begin{observation}[Osmotrophs evert the lumen]
\label{tub:obs:evert}
Observation~\ref{tub:obs:lumen} said the lumen is boundary namespace, owned but
open. An osmotroph runs the same converting chain without the ownership: the
substrate \emph{is} the lumen, and the enzyme sits at the boundary rather than
folded inside it.
\end{observation}

Observation~\ref{tub:obs:evert} is what \S\ref{tub:sec:organs} was reaching for
and could not yet say. There, letting a boundary namespace meant enclosure plus
tenancy. Here it is tenancy without enclosure. The fungus is the tube
construction turned inside out, which is why the mycorrhizal network appeared in
\chEng as the external instance of exactly the profile
Observation~\ref{tub:obs:gut} draws internally --- and a leaf and a root system
are the same purchase made by a non-phagotroph that is not an osmotroph either.

The cost structure is then legible rather than anomalous. Internalizing the
chain buys exclusivity over a partially extracted harvest, which is
Conjecture~\ref{tub:conj:intern}; externalizing it surrenders exclusivity and in
exchange never has to transport residue. An osmotroph is an agent that has taken
the other side of that trade, and it should be vulnerable exactly where the
trade is worst --- to theft of extracellular product by neighbors.

\subsection{A prior statement, and what is left to claim}

Honesty about priority. The functional claim is already in the literature: the
same review that catalogues the losses states in passing that a one-way gut
processes food more efficiently and permits uptake while the animal is still
digesting \cite{hejnol2015}.

What is claimed here is therefore not the observation but the derivation --- that
the two openings follow from concurrency together with ownership, and that the
threshold is a computable function of $\Lambda$ and $\mu$ rather than a
qualitative preference. Whether the derivation goes through is
Conjecture~\ref{tub:conj:intern} and the formal work below.

\section{What would have to be written down}
\label{tub:sec:formal}

\paragraph{The two-phase harvest as a term.} An exclusive acquisition rendezvous
on a boundary channel, in parallel with a persistent conversion chain on owned
channels, terminating in a send on the expulsion channel. The \texttt{Forage} of
the worked ecosystem in \chSci and the chain of \chEng supply most of it.

\paragraph{The tube as a namespace shape.} $\Chn_{\partial}(B)$ factors as
$\Chn_{\mathrm{in}}^{\partial} \sqcup \Chn_{\mathrm{out}}^{\partial}$ with a
directed redex pathway from intake to expulsion and no return path. Then
\emph{tube} is a formula in the generated logic rather than a designation, which
is the move already made for individuality itself in
Definition~\ref{eng:def:indiv}.

\paragraph{The occupancy ratio as a term-level quantity.} $\Lambda$ should be
recoverable from the harvest term itself: $\tau_{\mathrm{proc}}$ as the metered
cost of the conversion chain, $\tau_{\mathrm{enc}}$ as the expected waiting time
on the intake channel under the ambient offer rate. Then
Proposition~\ref{tub:prop:criterion} becomes an inequality between two
quantities the framework already meters, and the threshold is derived rather
than posited. This is the single most valuable item on the list.

\paragraph{Separability as a namespace condition.} $\mu$ is not a physical
parameter but a statement about whether residue and intake can share a name
without rendezvous --- that is, whether the sorts involved make the crossed
redex impossible. Gas exchange has $\mu = 1$ because the spent medium and the
resource are the same sort and the extraction happens at the wall rather than at
the aperture. Stated as a condition on the sort discipline,
Proposition~\ref{tub:prop:criterion} becomes entirely internal.

\paragraph{Whether the absence of a return path is forced.} It is not.
Rumination and caecotrophy are return paths, and they occur exactly where one
pass leaves too much unextracted: lagomorphs and hystricomorph rodents run a
colonic separation mechanism that diverts fine particles and microbes back to
the caecum and then re-ingest the product \cite{hirakawa2001}. This should be
stated as a condition on residual yield rather than suppressed, since a
proposition that forbids observed behavior is worse than one that predicts it.

\paragraph{The physical topology as a shadow.} A through-gut makes a body a
genus-one surface. That is real, but it is not the content. The content is the
channel topology; the surface is what an embodiment of that topology in
three-space looks like. Saying so explicitly keeps the argument from appearing
to run from anatomy to mathematics, which is the direction it does not run.

\section{Open questions}
\label{tub:sec:open}

\begin{question}
Is the tube forced or merely favored? Proposition~\ref{sci:prop:composite}
predicts mixotrophs, and predicts that their tubes should be facultative. Is that
borne out?
\end{question}

\begin{question}
Does the surface-area argument connect to the critical radius of
Corollary~\ref{eng:recovering-r}? A tube buys boundary area at fixed volume;
$r^\ast$ is a critical radius. There is plausibly one calculation here rather
than two.
\end{question}

\begin{question}
Where does circulation enter? Beyond a certain body size the tube's yield must
be distributed, which is a second internal medium with a different profile ---
relay rather than converting --- and a different topology, branching and
returning. If the tube follows from the access profile, does the vascular tree
follow from the tube together with the allometric constraint of \cite{west}?
\end{question}

\begin{question}
What is the assay analogue? Tasting is an assay run at the intake channel,
deciding whether to admit; vomiting is a refutation acted on after admission.
The trichotomy of Proposition~\ref{sci:prop:trichotomy} and budget-relative
refutation should both apply.
\end{question}

\begin{question}
Does the tube give a new grading dimension, or is it already covered by the
vector over typed namespaces? The suspicion is that the tube is a fact about
$\Chn$, which would be an independent argument for the fifth component of
\S\ref{eng:sec:indiv}.
\end{question}

\begin{question}
Does the osmotroph / phagotroph split have a framework-native characterization?
Observation~\ref{tub:obs:evert} says the osmotroph runs the converting chain
outside the ownership boundary. If that is right, the split should fall out of
where the chain sits relative to $\Chn_{\partial}$, and osmotrophy should carry
a predictable exposure to theft that phagotrophy does not.
\end{question}

\begin{question}
Is there a second critical point at high $\Lambda$? Ruminants and hindgut
fermenters run $\tau_{\mathrm{proc}}$ far above $\tau_{\mathrm{enc}}$ and answer
with buffering --- a sequence of chambers, and in the ruminant a deliberate
return path. Buffering is the standard answer to a throughput mismatch, and the
question is whether the framework predicts the chamber count.
\end{question}

\section{What is claimed}
\label{tub:sec:claimed}

Claimed: that the residue of a harvest is a consequence of prey being
computations rather than an added postulate
(Observation~\ref{tub:obs:refuse}); that a lumen is boundary namespace and
therefore not interior (Observation~\ref{tub:obs:lumen}); that a digestive organ
is a community tenanted in that namespace, which is why a market can operate
inside a firm without contradicting the firm (Observation~\ref{tub:obs:let});
that an osmotroph is the same construction everted, tenancy without enclosure
(Observation~\ref{tub:obs:evert}); and that intake and expulsion on one name
force an alternation (Proposition~\ref{tub:prop:twoends}).

Claimed with a different kind of support: the tube criterion,
Proposition~\ref{tub:prop:criterion}. It is not derived. What stands behind it is
the test set of \S\ref{tub:sec:testset} --- seven cases proposed as
counterexamples, of which one was factually wrong, four came out as the
occupancy ratio predicts, and two corrected the scope of the claim in a way that
made it cover cestodes, which no formulation in terms of heterotrophy could have
done. That is the support a criterion can have before the derivation exists, and
the honest description of it is that the criterion has survived the first attempt
to kill it.

Not claimed: the six conjectures, and in particular not
Conjecture~\ref{tub:conj:intern}, on which any derivation of the criterion
depends.

What the chapter does claim, and would defend, is a change in what kind of thing
an anatomy is. Part~\ref{part:ecology} has spent its length arguing that a
learner is a term among terms, paying; that the unit of learning moves; and that
where one individual ends is a question with an answer in the framework rather
than a matter of convention. If that is right, then the shape of an animal is not
a separate subject that biology happens to study alongside cognition. It is
another entry on the same bill. A learner that must open its food pays for
inquiry, and pays again for somewhere to put the opening --- and the second
payment, made in surface rather than in tokens, is the reason that almost
everything which thinks is also, in the end, a tube.
